<%= heading2 %>{Writing in Markdown}\label{sec:markdown}

<%# latexindent indents the \ifluatex branch, except for the verbatim Markdown body and \end{markdown} -%>
<% const mdIndent = (latexcompiler == "both") ? '  ' : ''; -%>
<% if (latexcompiler == "both") { -%>
\ifluatex
<% } -%>
<%= mdIndent %>If you are new to \LaTeX{}, you can start by writing your text in Markdown using the \href{https://ctan.org/pkg/markdown}{markdown} package.
<%= mdIndent %>\LaTeX{} then takes care of the typesetting.
<%= mdIndent %>Markdown and \LaTeX{} can be mixed:
<%= mdIndent %><%- Cref %>{tbl:markdown} is referenced from within the Markdown text below.
<%= mdIndent %>A longer part, such as a whole <% if (heading1 == "\\chapter") { %>chapter<% } else { %>section<% } %>, can be placed in a separate file included by \verb|\markdownInput{introduction.md}|.
<%= mdIndent %>The \href{https://github.com/latextemplates/markdown-latex-quickstart}{markdown-latex-quickstart} shows a complete document written in Markdown.

<%= mdIndent %><%- bexample %>
<%= mdIndent %>\begin{markdown}
Text can be *emphasized* or **bold**, and `code` is set in a monospaced font.
A citation looks like this [@mwe], and here is a footnote.^[This is an inline footnote.]

- first point
- second point
  - sub-point

1) firstly
2) secondly

| Right | Left | Centered |
|------:|:-----|:--------:|
|    12 | 12   |    12    |
|   123 | 123  |   123    |

: A table written in Markdown {#tbl:markdown}

<%- Cref %>{tbl:markdown} shows that tables get a caption and a label.
\end{markdown}
<%= mdIndent %><%- eexample %>
<% if (latexcompiler == "both") { -%>
\else
  This example requires Lua\LaTeX{}.
  Compile the document with Lua\LaTeX{} to see it.
\fi
<% } -%>
